% ==================================================
% CHAPTER 9: CONCLUSION
% ==================================================

\chapter{Conclusion}
\label{ch:conclusion}

\section{Summary}

The Hostel Room Swapping Service turns a problem students try to solve
one conversation at a time into a single computation over the whole
hostel. Because the only rooms in play are the ones the participants
already hold, every solution is a permutation of students over bed
slots, and its cycles are the swap chains. The engine scores rooms
against weighted hard and soft preferences, finds the best assignment
with the Hungarian algorithm and its alternatives with Murty's
algorithm, keeps only chains in which nobody loses, and offers each
student their own best chains, capped at eight members. Around the
engine, the project built a round workflow and an approval state
machine with expiry, a FastAPI backend with session authentication, a
PostgreSQL schema whose key invariants are enforced by the database, and
a React interface for students and wardens.

\section{Achievements}

\begin{enumerate}
  \item \textbf{Large gains with nobody worse off}: on the frozen
  30-student cohort one round raises mean satisfaction from 11\% to
  77\%, with 20 students in their exact requested room and zero worse
  off.
  \item \textbf{Executable, not only optimal}: Murty's alternatives cut
  the longest chain needing joint approval from 28 students to 11 at no
  cost in satisfaction.
  \item \textbf{A measured design choice}: a ten-cohort, 300-student
  benchmark chose the chain cap of eight, which reaches 58\% of students
  with chains that need far fewer approvals than uncapped ones.
  \item \textbf{No student trapped}: one refusal frees the whole chain,
  silence expires after three days, and a finished swap frees its
  members, enforced both in code and by a database index.
  \item \textbf{An auditable record}: allocation history that no client
  can rewrite.
  \item \textbf{A faster search}: Murty's subproblems now shrink as
  edges are fixed, making the search 4--13 times faster with identical
  results.
  \item \textbf{Engineering discipline}: a framework-free domain, 152
  automated tests run in CI on three Python versions, numbered
  migrations, and weekly journals of problems, causes and fixes.
\end{enumerate}

\section{Limitations}

\begin{itemize}
  \item The engine has been evaluated only on synthetic cohorts, not on
  real registrations for a live hostel.
  \item With the cap of eight, about 42\% of students in a 300-student
  cohort get no offer. A long cycle cannot be split, so these matches
  are lost, not delayed.
  \item The live search budget ($K = 50$) is smaller than the one
  benchmarked ($K = 200$), so live reach on large cohorts is lower than
  the benchmark figures. The faster search makes a larger $K$
  affordable, but the live default has not been raised yet.
  \item Shared rooms and roommate requests are supported and unit-tested
  but not part of the evaluation cohorts.
  \item The login endpoint has no rate limit.
  \item The service is deployed but has not yet run a round with real
  students.
\end{itemize}

\section{Future Work}

\begin{enumerate}
  \item Run a pilot round with real registrations for one hostel.
  \item Rate-limit login attempts.
  \item Raise the live $K$ now that the search is faster, after
  measuring how reach on large cohorts grows with it.
  \item Warm-start Murty's subproblems from the parent's dual
  potentials, bringing the search towards $O(Kn^3)$.
  \item Notify students by email when a chain that includes them is
  offered, so fewer proposals expire unanswered.
\end{enumerate}

\section{Final Remarks}

The central idea of the project held up: a hostel's swap requests can
be solved as one closed assignment problem, and the chains fall out of
its cycle structure instead of being searched for. The measurements
showed that the optimal answer is not the useful one. The useful answer
is the one a group of real students can agree on, which is why Murty's
alternatives, the Pareto filter, the chain cap and the expiry rules
matter as much as the Hungarian solve itself. The remaining work is
mostly about meeting real students and real data, not about redesigning
the approach.
